\chapter{From the equations to the code}\label{ch:equations-to-code}

This manual describes the physics; the user and developer manual describes how to run and
extend the code. \Cref{tab:concordance} connects the two: for each concept it identifies the equation
or section of this manual, the corresponding code interfaces, and the user-manual page that explains
their use. The complete lists of options, result fields, command-line arguments, observables
and probes are generated from the code into the reference pages of the user manual, and the
error messages with their remedies are collected on its troubleshooting page:
\begin{itemize}
\item \usermanual{docs/index.md};
\item \usermanual{docs/reference/options.md}, and the neighbouring reference pages;
\item \usermanual{docs/troubleshooting.md}.
\end{itemize}

% Generated by tools/gen_docs.py from docs/concordance.toml; do not edit.
\begin{small}\setlength{\tabcolsep}{4pt}
\begin{longtable}{>{\raggedright\arraybackslash}p{3.3cm}l>{\raggedright\arraybackslash}p{4.9cm}>{\raggedright\arraybackslash}p{3.6cm}}
\caption{Concepts of this manual, their equations or sections, the corresponding names in the code, and the page of the user manual that explains their use.}\label{tab:concordance}\\
\toprule Concept & Here & In the code & User manual \\ \midrule\endfirsthead
\toprule Concept & Here & In the code & User manual \\ \midrule\endhead
\bottomrule\endfoot
Superconducting Anderson impurity and reference junction & \cref{eq:reference} & \code{reference_model}, \code{Impurity.}\allowbreak \code{anderson}, \code{make_model} & Physical models and conventions; tutorial 01 \\
Discretized superconducting reservoirs & \cref{sec:bath} & \code{cosh_grid}, \code{fit_surrogate}, \code{chain_expansion}, \code{DiscreteBath} & Choosing and refining superconducting reservoirs; tutorial 05 \\
Symmetry sectors & \cref{sec:symmetries} & \code{Sector} & Physical models and conventions; tutorial 01 \\
Canonical coordinates & \cref{sec:coordinates} & \code{Hamiltonian.}\allowbreak \code{physical_operator}, \code{Hamiltonian.}\allowbreak \code{coordinate_fingerprint} & Physical models and conventions \\
Quasiparticle expansion & \cref{eq:qp-projection} & \code{qp_solver.}\allowbreak \code{solve}, \code{qp_solver.}\allowbreak \code{SolverOptions} & QP calculations and observable-specific convergence; tutorial 00 \\
Density-matrix renormalization group & \cref{sec:dmrg} & \code{dmrg_solver.}\allowbreak \code{solve}, \code{dmrg_solver.}\allowbreak \code{SolverOptions} & DMRG states and bond-dimension refinement; tutorial 11 \\
Checked calculations and saved results & \cref{ch:practical} & \code{solve}, \code{save_result}, \code{load_result} & Parameter scans, saved results and reproducible analysis; tutorial 01 \\
Supercurrent & \cref{eq:current} & \code{phase_derivative} & Physical models and conventions; tutorial 01 \\
Transition frequencies and static matrix elements & \cref{eq:transitions} & \code{Eigenstates.}\allowbreak \code{matrix_elements}, \code{Hamiltonian.}\allowbreak \code{probe} & Using the solver and saving results; tutorial 06 \\
Parameter drives and harmonics & \cref{eq:drive} & \code{ParameterDrive}, \code{Harmonic}, \code{PeriodicHamiltonian.}\allowbreak \code{from_parameters} & Floquet calculations; tutorial 06 \\
Floquet states and quasienergies & \cref{eq:floquet-state} & \code{FloquetStates.}\allowbreak \code{quasienergies}, \code{FloquetStates.}\allowbreak \code{modes} & Floquet calculations; tutorial 02 \\
Mean energy & \cref{eq:mean-energy} & \code{FloquetStates.}\allowbreak \code{mean_energies} & Floquet calculations; tutorial 02 \\
Sambe matrix and Floquet cutoff & \cref{eq:sambe-matrix} & \code{solve_floquet}, \code{FloquetOptions.}\allowbreak \code{floquet_cutoff} & Floquet calculations; tutorial 05 \\
Static basis of size K & \cref{sec:sambe-method} & \code{FloquetOptions.}\allowbreak \code{basis}, \code{FloquetOptions.}\allowbreak \code{basis_states} & Floquet calculations; tutorial 05 \\
Replicas, branches and folded quasienergies & \cref{sec:sambe-method} & \code{FloquetOptions.}\allowbreak \code{selection}, \code{FloquetOptions.}\allowbreak \code{branch}, \code{FloquetStates.}\allowbreak \code{folded} & Floquet calculations \\
Residual bound & \cref{eq:residual-bound} & \code{FloquetStates.}\allowbreak \code{residuals}, \code{FloquetStates.}\allowbreak \code{fock_residuals} & Floquet calculations; tutorial 05 \\
Gauge freedom and c-number terms & \cref{sec:gauge} & \code{FloquetOptions.}\allowbreak \code{drop_scalar_harmonics} & Floquet calculations; tutorial 02 \\
Fixed canonical coordinates & \cref{sec:fixed-coordinates} & \code{PeriodicHamiltonian.}\allowbreak \code{coordinate_fingerprint} & Floquet calculations \\
One-period propagator & \cref{eq:cfm4} & \code{floquet_propagator}, \code{FloquetOptions.}\allowbreak \code{steps} & Floquet calculations; tutorial 05 \\
High-frequency expansion & \cref{eq:van-vleck} & \code{PeriodicHamiltonian.}\allowbreak \code{van_vleck} & Floquet calculations; tutorial 04 \\
Rotating-wave model & \cref{eq:rwa} & \code{rwa_model} & Floquet calculations; tutorial 03 \\
Floquet matrix elements & \cref{eq:floquet-element} & \code{FloquetStates.}\allowbreak \code{matrix_elements}, \code{FloquetTransitions} & Floquet calculations; tutorial 08 \\
Floquet golden rule and lines & \cref{eq:floquet-golden} & \code{FloquetTransitions.}\allowbreak \code{lines}, \code{SpectralLines.}\allowbreak \code{spectrum}, \code{line_shape} & Floquet calculations; tutorial 08 \\
Sum rule & \cref{eq:sum-rule} & \code{FloquetTransitions.}\allowbreak \code{completeness} & Floquet calculations; tutorial 02 \\
\end{longtable}
\end{small}


\section{Reproducing the figures and this manual}\label{sec:reproduce}

\begin{lstlisting}[language=bash]
cd docs/manual
make data      # recompute data/*.json with the solver (tens of minutes)
make quoted    # numbers quoted in the text, from data/*.json
make figures   # render figures/*.pdf from data/*.json (requires matplotlib)
make pdf       # quoted values and figures, then latexmk (pdflatex + biber)
\end{lstlisting}

Each record in \code{data/} stores the script name, the code version and git commit, and
the numerical environment. Figures and quoted numerical values are generated from the saved
data without rerunning the solver, so the PDF can be rebuilt quickly. The figure
scripts write the quoted values to \code{generated/quoted.tex}. Recomputing the data
therefore updates the values in the text, and changes can be reviewed in the generated-file diff.
